\section{Kết luận}

Trong bài báo này, chúng tôi đã giới thiệu một framework mới dựa trên ViT cho ZSL nhằm tăng cường sự liên kết thị giác-ngữ nghĩa bằng cách giải quyết vấn đề sai lệch ngữ nghĩa ở không gian đầu vào. Phương pháp của chúng tôi bao gồm hai thành phần chính: lựa chọn bản vá tự giám sát, tận dụng các điểm chú ý tổng hợp để xác định và loại bỏ các bản vá không liên quan đến ngữ nghĩa, và Bối cảnh hóa Ngữ nghĩa, thay thế các bản vá này bằng các nhúng có thể học được và nhận biết thuộc tính. Các thử nghiệm sâu rộng trên các bộ dữ liệu ZSL tiêu chuẩn đã chứng minh rằng phương pháp của chúng tôi đạt được kết quả hiệu suất hiện đại. Công trình của chúng tôi nhấn mạnh tầm quan trọng của việc tinh chỉnh ngữ nghĩa ở giai đoạn đầu trong ZSL và đề xuất các hướng đi mới để tích hợp các tiên nghiệm ngữ nghĩa vào các mô hình thị giác dựa trên transformer. Công việc trong tương lai sẽ khám phá các chiến lược thích ứng để lựa chọn và bối cảnh hóa bản vá, cho phép mô hình điều chỉnh linh hoạt dựa trên các đặc điểm của bộ dữ liệu.
